% ---------------------------------------------------------------------
% Discussion, Limitations, and Ethics (~1 page).
% ---------------------------------------------------------------------

\subsection{What the Results Mean for Practice}

Three practical lessons emerge from Section~\ref{sec:results}. First, \emph{in-domain
leaderboards should no longer drive countermeasure selection}: every system in
Table~\ref{tab:main} sits below 1.1\% EER in-domain while spanning a factor of five in
zero-shot error. Procurement or deployment decisions based on ASVspoof~2019 numbers
alone are uninformative. Second, \emph{robustness is a training-data property, not an
architectural one}: the same architecture spans 4.6--9.8\% EER under 64~kbps MP3
depending only on whether the augmentation curriculum was used
(Table~\ref{tab:robust}). Practitioners can retrofit most of this gain onto existing
detectors. Third, \emph{calibration must be validated under shift}: temperature scaling
fitted in-domain transferred acceptably for the one-class score but left the BCE
baselines over-confident on In-the-Wild (Table~\ref{tab:calibration}); a deployed
system should monitor calibration on traffic, not trust laboratory ECE.

\subsection{Why Not Just Use a Foundation Model?}

Our ALLM results (B8) bear on a question the field will face repeatedly as audio
foundation models improve: whether general-purpose models subsume specialized
countermeasures. Zero-shot Qwen2-Audio performs near chance across all seven corpora,
consistent with ALLM4ADD~\cite{gong2025allm4add}. The failure is informative rather
than surprising: instruction-tuned ALLMs are optimized to extract \emph{semantic}
content, and their audio tokenizers---often themselves neural codecs---discard the
phase and fine spectral residue that constitutes spoofing evidence. Indeed, Codecfake's
premise~\cite{xie2025codecfake} is that codec resynthesis \emph{is} a spoofing attack;
an ALLM that perceives audio through such a codec is structurally blind to the
evidence. We expect fine-tuned ALLMs~\cite{gong2025allm4add} to become competitive on
seen conditions, but the calibration and open-set properties that forensic use requires
remain, for now, the domain of purpose-built detectors. A promising hybrid is to use
AuralGuard-style scores as tool outputs that an ALLM can reason over when asked to
assess authenticity holistically.

\subsection{Limitations}

\textbf{Adaptive adversaries.} Our threat model covers unseen generators and channels,
not attackers who optimize against the deployed detector. ASVspoof~5's
surrogate-optimized adversarial attacks~\cite{wang2025asvspoof5journal} are included in
the zero-shot evaluation, but white-box adaptive attacks would require dedicated
defenses; we deliberately do not publish an anti-forensic recipe.
\textbf{Compute footprint.} The frozen WavLM-Large encoder dominates inference cost
(Table~\ref{tab:budget}). CPU real-time factor of $\approx$0.4 suffices for server-side
moderation but not for on-device screening; distillation into lightweight
back-ends~\cite{liu2025nes2net} is a natural next step.
\textbf{Curriculum coverage.} The augmentation chain approximates today's dominant
transmission paths; emerging channels (e.g., neural-codec voice calls) will require
curriculum updates, and our own results show robustness does not extrapolate to
distortions outside the curriculum family.
\textbf{Partial spoofs and other modalities.} We score whole utterances; localized
splicing~\cite{luong2024llamapartialspoof} and singing-voice
deepfakes~\cite{zang2024singfake} are out of scope, though the windowed inference of
Section~\ref{ssec:trainproc} provides a straightforward extension path.
\textbf{Language coverage.} MLAAD results (Table~\ref{tab:crosslingual}) show a
persistent gradient against low-resource language families inherited from the SSL
backbone's pre-training distribution; claims about universal applicability would be
premature.

\subsection{Ethical Considerations}

Audio deepfake detection is dual-use. AuralGuard is intended to reduce harm from
voice-cloning fraud and audio disinformation, but detectors can also lend false
authority: a miscalibrated ``97\% synthetic'' verdict can wrongly discredit genuine
recordings---the reverse harm of the one it prevents, and one documented in real
fact-checking practice~\cite{chandra2025deepfakeeval}. Our system therefore
(i)~outputs calibrated probabilities with explicit uncertainty rather than binary
verdicts, (ii)~exposes per-decision evidence (gate statistics and attribution
heatmaps), and (iii)~is documented, in both the paper and the released demo, as a
triage aid that must not be the sole basis for consequential decisions. All training
and evaluation corpora are public research datasets collected under their own consent
frameworks; no private voice data was scraped. We audit performance across languages
(Table~\ref{tab:crosslingual}) and report the disparities we find rather than a single
pooled number. Finally, we release code and weights to enable scrutiny and defensive
research, while withholding adaptive-attack tooling.
